\documentclass[11pt]{article}

% ---------- Packages ----------
\usepackage[margin=1in]{geometry}
\usepackage{amsmath, amssymb, amsthm}
\usepackage{enumitem}
\usepackage{fancyhdr}
\usepackage{lastpage}
\usepackage{graphicx}
\usepackage[colorlinks=true, linkcolor=blue, urlcolor=blue, citecolor=blue]{hyperref}

% ---------- Header / Footer ----------
\pagestyle{fancy}
\fancyhf{}
\lhead{STAT 20000: Elementary Statistics}
\rhead{Homework 3}
\cfoot{Page \thepage\ of \pageref{LastPage}}

% ---------- Custom Commands ----------
\newcommand{\course}{STAT 20000: Elementary Statistics}
\newcommand{\hwtitle}{Homework 3}
\newcommand{\duedate}{Due: 11:59 pm, April 17, 2026}
% \newcommand{\studentname}{Your Name}

% ---------- Question Environment ----------
% ---------- Question Environment with Points ----------
\newcounter{question}
\newenvironment{question}[2][]{
  \refstepcounter{question}
  \vspace{1em}
  \noindent\textbf{Question \thequestion\ (#2 points). #1}
}{\vspace{1em}}

% Parts (a), (b), (c)
\newenvironment{parts}{
  \begin{enumerate}[label=(\alph*)]
}{
  \end{enumerate}
}

% Optional solution environment (comment out to hide solutions)
\newenvironment{solution}{
  \par\noindent\textit{Solution.}
}{
  \vspace{1em}
}

% ---------- Document ----------
\begin{document}

\begin{center}
    {\Large \textbf{\course}}\\
    {\large \hwtitle}\\
    \duedate\\
    \vspace{0.5em}
    % \textbf{Name:} \studentname
\end{center}

\hrule
\vspace{1em}

% ---------- Questions ----------

\noindent All page, section, and exercise numbers below refer to the course text \textit{Statistics}, 4th edition, by D.Freedman, R. Pisani \& R. Purves (FPP).

\paragraph{Reading: }Chapters 5, 8 and 9 in FPP.

\paragraph{Q1 (3 points).} (Chapter 5, Review Exercise 6 on p.95)
Among applicants to one law school, the average LSAT score was about 169,
the SD was about 9, and the highest score was 178. Did the LSAT scores
follow the normal curve?

\paragraph{Q2 (8 points)} (Chapter 5, Review Exercise 7 on p.95) Among freshmen at a certain university, scores on the Math SAT followed the normal curve, with an average of 550 and an SD of 100. Fill in the blanks; explain briefly.

\begin{itemize}
    \item [(a)] A student who scored 400 on the Math SAT was at the \underline{~~~~~~~~}th percentile of the score distribution.
    \item [(b)] To be the 75th percentile of the distribution, a student needed a score of about \underline{~~~~~~~~} points on the Math SAT
\end{itemize}


\paragraph{Q3 (3 points).} (Chapter 9, Review Exercise 6 on p.154) A computer program prints out $r$ for the two data sets shown below. Is the program working correctly? Answer yes or no, and explain briefly.

\medskip

\noindent
\begin{minipage}{0.45\textwidth}
\begin{center}
(i)

\begin{tabular}{cc}
$x$ & $y$ \\ \hline
1 & 2 \\
2 & 1 \\
3 & 4 \\
4 & 3 \\
5 & 7 \\
6 & 5 \\
7 & 6 \\
\end{tabular}

\medskip
$r = 0.8214$
\end{center}
\end{minipage}
\hfill
\begin{minipage}{0.45\textwidth}
\begin{center}
(ii)

\begin{tabular}{cc}
$x$ & $y$ \\ \hline
1 & 5 \\
2 & 4 \\
3 & 7 \\
4 & 6 \\
5 & 10 \\
6 & 8 \\
7 & 9 \\
\end{tabular}

\medskip
$r = 0.7619$
\end{center}
\end{minipage}


\paragraph{Q4 (3 points)} (Chapter 8, Review Exercise 3 on p.135) Suppose men always married women who were exactly 8\% shorter. What would the correlation between their heights be?

\paragraph{Q5 (6 points)} (Chapter 5, Review Exercise 8 on p.95) True or false? Explain briefly.
\begin{itemize}
    \item[(a)] If you add 7 to each entry on a list, that adds 7 to the average.
    \item[(b)] If you add 7 to each entry on a list, that adds 7 to the SD.
    \item[(c)] If you double each entry on a list, that doubles the average.
    \item[(d)] If you double each entry on a list, that doubles the SD.
    \item[(e)] If you change the sign of each entry on a list, that changes the sign of the average.
    \item[(f)] If you change the sign of each entry on a list, that changes the sign of the SD.
\end{itemize}

\paragraph{Q6 (2 points)} (Chapter 9, Review Exercise 11 on p.156) The Educational Testing Service computed the average Verbal SAT score for each state, as well as the average
Math SAT score for each state. (D.C. counts as a state.) The correlation between these $51$ pairs of averages was $0.97$. Would the correlation
between the Math SAT and the Verbal SAT—computed from the data on all
the individuals who took the tests—be larger than $0.97$, about $0.97$, or less
than $0.97$? Explain briefly.

\end{document}